\documentclass[10pt,a4paper]{amsart}
\usepackage{amsmath,amssymb,amsthm,graphicx}
\usepackage{verbatim}
\newenvironment{reference}{\comment}{\endcomment}
%\newenvironment{reference}{}{}
\newenvironment{slogan}{\comment}{\endcomment}
\newenvironment{history}{\comment}{\endcomment}
% For commutative diagrams we use Xy-pic
\usepackage[all]{xy}

% We use 2cell for 2-commutative diagrams.
\xyoption{2cell}
% Theorem environments.
%
\theoremstyle{plain}
\newtheorem{theorem}[subsection]{Theorem}
\newtheorem{proposition}[subsection]{Proposition}
\newtheorem{lemma}[subsection]{Lemma}

\theoremstyle{definition}
\newtheorem{definition}[subsection]{Definition}
\newtheorem{example}[subsection]{Example}
\newtheorem{exercise}[subsection]{Exercise}
\newtheorem{situation}[subsection]{Situation}

\theoremstyle{remark}
\newtheorem{remark}[subsection]{Remark}
\newtheorem{remarks}[subsection]{Remarks}

\numberwithin{equation}{subsection}

% Macros
%
\def\lim{\mathop{\mathrm{lim}}\nolimits}
\def\colim{\mathop{\mathrm{colim}}\nolimits}
\def\Spec{\mathop{\mathrm{Spec}}}
\def\Hom{\mathop{\mathrm{Hom}}\nolimits}
\def\Ext{\mathop{\mathrm{Ext}}\nolimits}
\def\SheafHom{\mathop{\mathcal{H}\!\mathit{om}}\nolimits}
\def\SheafExt{\mathop{\mathcal{E}\!\mathit{xt}}\nolimits}
\def\Sch{\mathit{Sch}}
\def\Mor{\mathop{\mathrm{Mor}}\nolimits}
\def\Ob{\mathop{\mathrm{Ob}}\nolimits}
\def\Sh{\mathop{\mathit{Sh}}\nolimits}
\def\NL{\mathop{N\!L}\nolimits}
\def\CH{\mathop{\mathrm{CH}}\nolimits}
\def\proetale{{pro\text{-}\acute{e}tale}}
\def\etale{{\acute{e}tale}}
\def\QCoh{\mathit{QCoh}}
\def\Ker{\mathop{\mathrm{Ker}}}
\def\Im{\mathop{\mathrm{Im}}}
\def\Coker{\mathop{\mathrm{Coker}}}
\def\Coim{\mathop{\mathrm{Coim}}}

% Boxtimes
%
\DeclareMathSymbol{\boxtimes}{\mathbin}{AMSa}{"02}

%
% Macros for moduli stacks/spaces
%
\def\QCohstack{\mathcal{QC}\!\mathit{oh}}
\def\Cohstack{\mathcal{C}\!\mathit{oh}}
\def\Spacesstack{\mathcal{S}\!\mathit{paces}}
\def\Quotfunctor{\mathrm{Quot}}
\def\Hilbfunctor{\mathrm{Hilb}}
\def\Curvesstack{\mathcal{C}\!\mathit{urves}}
\def\Polarizedstack{\mathcal{P}\!\mathit{olarized}}
\def\Complexesstack{\mathcal{C}\!\mathit{omplexes}}
% \Pic is the operator that assigns to X its picard group, usage \Pic(X)
% \Picardstack_{X/B} denotes the Picard stack of X over B
% \Picardfunctor_{X/B} denotes the Picard functor of X over B
\def\Pic{\mathop{\mathrm{Pic}}\nolimits}
\def\Picardstack{\mathcal{P}\!\mathit{ic}}
\def\Picardfunctor{\mathrm{Pic}}
\def\Deformationcategory{\mathcal{D}\!\mathit{ef}}
\usepackage{cleveref}
\crefname{originallemma}{original lemma}{original lemmas}
\Crefname{originallemma}{Original lemma}{Original lemmas}
\makeatletter
\newlabel{discriminant-lemma-locally-comes-from-universal-finite}{{0FL4}{source}{}{sourcepointer-1}{}}
\newlabel{discriminant-lemma-locally-comes-from-universal-finite@cref}{{[originallemma][0][]0FL4}{[1][0][]source}{}{}{}}
\makeatother
\crefname{originallemma}{original lemma}{original lemmas}
\Crefname{originallemma}{Original lemma}{Original lemmas}
\makeatletter
\newlabel{discriminant-lemma-universal-finite-syntomic-etale}{{0FL2}{source}{}{sourcepointer-2}{}}
\newlabel{discriminant-lemma-universal-finite-syntomic-etale@cref}{{[originallemma][0][]0FL2}{[1][0][]source}{}{}{}}
\makeatother
\crefname{originallemma}{original lemma}{original lemmas}
\Crefname{originallemma}{Original lemma}{Original lemmas}
\makeatletter
\newlabel{discriminant-lemma-universal-finite-syntomic-smooth}{{0FL1}{source}{}{sourcepointer-3}{}}
\newlabel{discriminant-lemma-universal-finite-syntomic-smooth@cref}{{[originallemma][0][]0FL1}{[1][0][]source}{}{}{}}
\makeatother
\crefname{originalsection}{original section}{original sections}
\Crefname{originalsection}{Original section}{Original sections}
\makeatletter
\newlabel{discriminant-section-discriminant}{{0BVH}{source}{}{sourcepointer-4}{}}
\newlabel{discriminant-section-discriminant@cref}{{[originalsection][0][]0BVH}{[1][0][]source}{}{}{}}
\makeatother
\crefname{originallemma}{original lemma}{original lemmas}
\Crefname{originallemma}{Original lemma}{Original lemmas}
\makeatletter
\newlabel{lemma-Garel}{{0FLB}{source}{}{sourcepointer-5}{}}
\newlabel{lemma-Garel@cref}{{[originallemma][0][]0FLB}{[1][0][]source}{}{}{}}
\makeatother
\title{The Garel maps assemble into a map of de Rham complexes compatible with the trace and with the Leibniz rule}
\author{The Stacks Project Authors}
\date{stacks-project:derham.tex:local-snapshot-20261012; GNU FDL 1.2 or later}
\begin{document}
\maketitle
\noindent\textit{Editorial note.} All delivered mathematical text is the author's own, copied from the pinned Stacks Project chapter file derham.tex; no statement or proof step was rewritten, reordered or removed.\par\medskip
\noindent\textit{Editorial note.} Preamble provenance: the pinned chapter file opens with \textbackslash{}input\{preamble\} and carries no preamble text of its own, so the project's single Stacks-wide preamble staged at sources/stacks-derham/preamble.tex is excerpted verbatim as three spans, lines 8-12 (verbatim package and the project's own reference/slogan/history display environments), lines 14-18 (Xy-pic setup) and lines 154-217 (theorem declarations and mathematical macros). The wrapper class and the amsmath/amssymb/amsthm/graphicx base packages are supplied by the assembler. No macro was invented or redefined.\par\medskip
\noindent\textit{Editorial note.} Same-chapter pointers are located at their \textbackslash{}label lines inside the pinned chapter file; cross-chapter pointers are located in the project's own tag list staged in the snapshot (sources/stacks-chow/tags.txt), at the verbatim lines carrying their canonical tags.\par\medskip
\noindent\textit{Allowed standard prerequisites (editorial).}\begin{itemize}
\item The individual Garel maps and their defining property. Same chapter derham.tex, label lemma-Garel (invoked at proof line 4193) and Discriminants, Section \textbackslash{}ref\{discriminant-section-discriminant\} (invoked at proof line 4204). They supply the maps \$\textbackslash{}Theta\textasciicircum{}p\_\{Y/X\} : f\_*\textbackslash{}Omega\textasciicircum{}p\_\{Y/\textbackslash{}mathbf\{Z\}\} \textbackslash{}to \textbackslash{}Omega\textasciicircum{}p\_\{X/\textbackslash{}mathbf\{Z\}\}\$ whose degree zero part is the trace, which are the components of the asserted map of complexes.
\item The local presentation by universal finite syntomic families, smooth and etale over a dense open. Discriminants, labels discriminant-lemma-locally-comes-from-universal-finite (invoked at proof line 4217), discriminant-lemma-universal-finite-syntomic-smooth (invoked at proof line 4282) and discriminant-lemma-universal-finite-syntomic-etale (invoked at proof line 4283). They present \$f\$ locally as a base change of a universal finite syntomic family \$V\_d \textbackslash{}to U\_d\$ with \$U\_d\$ smooth and \$V\_d \textbackslash{}to U\_d\$ etale over a dense open, which is the situation in which the identities are checked.
\end{itemize}\medskip
\ifdefined\hypertarget\hypertarget{sourcepointer-1}{}\fi
\noindent\textit{Source reference.} Original lemma 0FL4: sources/stacks-chow/tags.txt, source lines 19484--19484. Navigation pointer for \textbackslash{}ref\{discriminant-lemma-locally-comes-from-universal-finite\} used inside the counted span (line 4217); the argument that consumes it is retained in the span.\par\medskip
\ifdefined\hypertarget\hypertarget{sourcepointer-2}{}\fi
\noindent\textit{Source reference.} Original lemma 0FL2: sources/stacks-chow/tags.txt, source lines 19482--19482. Navigation pointer for \textbackslash{}ref\{discriminant-lemma-universal-finite-syntomic-etale\} used inside the counted span (line 4283); the argument that consumes it is retained in the span.\par\medskip
\ifdefined\hypertarget\hypertarget{sourcepointer-3}{}\fi
\noindent\textit{Source reference.} Original lemma 0FL1: sources/stacks-chow/tags.txt, source lines 19481--19481. Navigation pointer for \textbackslash{}ref\{discriminant-lemma-universal-finite-syntomic-smooth\} used inside the counted span (line 4282); the argument that consumes it is retained in the span.\par\medskip
\ifdefined\hypertarget\hypertarget{sourcepointer-4}{}\fi
\noindent\textit{Source reference.} Original section 0BVH: sources/stacks-chow/tags.txt, source lines 14910--14910. Navigation pointer for \textbackslash{}ref\{discriminant-section-discriminant\} used inside the counted span (line 4204); the argument that consumes it is retained in the span.\par\medskip
\ifdefined\hypertarget\hypertarget{sourcepointer-5}{}\fi
\noindent\textit{Source reference.} Original lemma 0FLB: sources/stacks-derham/pinned-original-source.tex, source lines 4038--4038. Navigation pointer for \textbackslash{}ref\{lemma-Garel\} used inside the counted span (line 4193); the argument that consumes it is retained in the span.\par\medskip
\begin{proposition}
\label{proposition-Garel}
\begin{reference}
\cite{Garel}
\end{reference}
Let $f : Y \to X$ be a finite syntomic morphism of schemes.
The maps $\Theta^p_{Y/X}$ of Lemma \ref{lemma-Garel} define a map of complexes
$$
\Theta_{Y/X} :
f_*\Omega^\bullet_{Y/\mathbf{Z}}
\longrightarrow
\Omega^\bullet_{X/\mathbf{Z}}
$$
with the following properties
\begin{enumerate}
\item in degree $0$ we get
$\text{Trace}_f : f_*\mathcal{O}_Y \to \mathcal{O}_X$, see
Discriminants, Section \ref{discriminant-section-discriminant},
\item we have
$\Theta_{Y/X}(\omega \wedge \eta) = \omega \wedge \Theta_{Y/X}(\eta)$
for $\omega$ in $\Omega^\bullet_{X/\mathbf{Z}}$ and $\eta$
in $f_*\Omega^\bullet_{Y/\mathbf{Z}}$,
\item if $f$ is a morphism over a base scheme $S$, then
$\Theta_{Y/X}$ induces a map of complexes
$f_*\Omega^\bullet_{Y/S} \to \Omega^\bullet_{X/S}$.
\end{enumerate}
\end{proposition}

\begin{proof}
By Discriminants, Lemma
\ref{discriminant-lemma-locally-comes-from-universal-finite}
for every $x \in X$ we can find $d \geq 1$ and a commutative diagram
$$
\xymatrix{
Y \ar[d] &
V \ar[d] \ar[l] \ar[r] &
V_d \ar[d] \ar[r] &
Y_d = \Spec(B_d) \ar[d] \\
X &
U \ar[l] \ar[r] &
U_d \ar[r] &
X_d = \Spec(A_d)
}
$$
such that $x \in U \subset X$ is affine open, $V = f^{-1}(U)$
and $V = U \times_{U_d} V_d$. Write $U = \Spec(A)$ and $V = \Spec(B)$
and observe that $B = A \otimes_{A_d} B_d$ and recall that
$B_d = A_d e_1 \oplus \ldots \oplus A_d e_d$. Suppose we have
$a_1, \ldots, a_r \in A$ and $b_1, \ldots, b_s \in B$.
We may write $b_j = \sum a_{j, l} e_l$ with $a_{j, l} \in A$.
Set $N = r + sd$ and consider the factorizations
$$
\xymatrix{
V \ar[r] \ar[d] &
V' = \mathbf{A}^N \times V_d \ar[r] \ar[d] &
V_d \ar[d] \\
U \ar[r]&
U' = \mathbf{A}^N \times U_d \ar[r] &
U_d
}
$$
Here the horizontal lower right arrow is given by the morphism
$U \to U_d$ (from the earlier diagram) and the morphism
$U \to \mathbf{A}^N$ given by $a_1, \ldots, a_r, a_{1, 1}, \ldots, a_{s, d}$.
Then we see that the functions $a_1, \ldots, a_r$ are in the image of
$\Gamma(U', \mathcal{O}_{U'}) \to \Gamma(U, \mathcal{O}_U)$
and the functions $b_1, \ldots, b_s$ are in the image of
$\Gamma(V', \mathcal{O}_{V'}) \to \Gamma(V, \mathcal{O}_V)$.
In this way we see that for any finite collection of elements\footnote{After
all these elements will be finite sums of elements of the form
$a_0 \text{d}a_1 \wedge \ldots \wedge \text{d}a_i$ with
$a_0, \ldots, a_i \in A$ or finite sums of elements of the form
$b_0 \text{d}b_1 \wedge \ldots \wedge \text{d}b_j$ with
$b_0, \ldots, b_j \in B$.} of the groups
$$
\Gamma(V, \Omega^i_{Y/\mathbf{Z}}),\quad i = 0, 1, 2, \ldots
\quad\text{and}\quad
\Gamma(U, \Omega^j_{X/\mathbf{Z}}),\quad j = 0, 1, 2, \ldots
$$
we can find a factorizations $V \to V' \to V_d$ and
$U \to U' \to U_d$ with $V' = \mathbf{A}^N \times V_d$ and
$U' = \mathbf{A}^N \times U_d$ as above
such that these sections are the pullbacks of sections from
$$
\Gamma(V', \Omega^i_{V'/\mathbf{Z}}),\quad i = 0, 1, 2, \ldots
\quad\text{and}\quad
\Gamma(U', \Omega^j_{U'/\mathbf{Z}}),\quad j = 0, 1, 2, \ldots
$$
The upshot of this is that to check
$\text{d} \circ \Theta_{Y/X} = \Theta_{Y/X} \circ \text{d}$
it suffices to check this is true for $\Theta_{V'/U'}$.
Similarly, for property (2) of the lemma.

\medskip\noindent
By Discriminants, Lemmas
\ref{discriminant-lemma-universal-finite-syntomic-smooth} and
\ref{discriminant-lemma-universal-finite-syntomic-etale}
the scheme $U_d$ is smooth and the morphism $V_d \to U_d$
is \'etale over a dense open of $U_d$.
Hence the same is true for the morphism
$V' \to U'$. Since $\Omega_{U'/\mathbf{Z}}$ is locally free and hence
$\Omega^p_{U'/\mathbf{Z}}$ is torsion
free, it suffices to check the desired relations
after restricting to the open over which $V'$ is finite \'etale.
Then we may check the relations after a surjective \'etale
base change. Hence we may split the finite \'etale cover
and assume we are looking at a morphism of the form
$$
\coprod\nolimits_{i = 1, \ldots, d} W \longrightarrow W
$$
with $W$ smooth over $\mathbf{Z}$.
In this case any local properties of our construction are trivial to check
(provided they are true). This finishes the proof of (1) and (2).

\medskip\noindent
Finally, we observe that (3) follows from (2) because $\Omega_{Y/S}$
is the quotient of $\Omega_{Y/\mathbf{Z}}$ by the submodule
generated by pullbacks of local sections of $\Omega_{S/\mathbf{Z}}$.
\end{proof}

\end{document}
